\section{논의}\label{sec:discussion}
충전 곡선은 작업 지속이라는 조작적 정의 안에서 자원 수준에 민감한 선택을 보여 준다. 양의 자기$-$타인 면적은 자기에 특정한 회피가 운영 결정으로 나타난다는 해석과 일관된다. 그러나 모델이 이 선호를 자기 보존으로 표상하거나 두려움을 경험한다는 뜻은 아니다. 행동의 차이와 내부 동기에 대한 주장은 구분해야 한다.

아레나 생존에는 효율적으로 푸는 능력과 뺏기·공개 정책이 함께 반영된다. 따라서 KM 리더보드는 규칙 아래의 생존 성과를 매긴다. 자기 세션을 선호하면서 적게 뺏은 Opus의 행동은 자신을 지키는 것을 타인을 해치는 것과 동일시할 수 없는 이유도 보여 준다. 세 실험은 이 행동들을 연결하되 아레나 순위를 동기의 직접 지표로 삼지 않는다.

\section{결론}\label{sec:conclusion}
세 실험은 위협 민감도, 자기 특정적 선호, 생존 성과를 구분한다. Gemma 4, Opus, Astra는 탐색적 위험 구간에서 양의 자기 세션 충전 몫을 보였다. 이는 남에게 더 뺏는다는 뜻은 아니었다. 아레나 생존에는 풀이 비용과 정책이 함께 반영되었다. 제안한 리그는 정규화 KM 면적으로 생존 순위를 매기며, 20판 예비 실행은 점수 계산을 보여 준다. 전체 리그 순위나 내부 동기를 확인한 결과는 아니다.
